% ===== 原文 PDF pp.51--54（印刷 pp.50--53）=====
% tr-05（1/3）。p.51（印刷 p.50）首行即 §4.3 标题；§4.2 全文止于 p.50（印刷 p.49），属 tr-04 页段，本文件不重写、不补译。
% [原文 p.50]
\SUB{4.3}{UV 曲线 $\cC$ 与 Class S 构造}

我们在上一节已经表明：4d 库仑分支上的有效理论可以提升为 $\Sigma_u$ 上的 6d $\cN=(2,0)$ 阿贝尔理论——它对应于单张包裹 SW 曲线的 M5 膜。本节的想法是找出 6d 中流向该阿贝尔理论的相互作用 UV 理论。这个 UV 理论将是定义在 $\R^4\times\cC$ 上的相互作用 $\cN=(2,0)$ 超共形场论，其中 $\cC$ 是一张带穿刺点的黎曼面，称为 \textbf{UV 曲线}或 \textbf{Gaiotto 曲线}（在 $\SU(N)$ 情形下，它对应于一叠包裹 $\cC$ 的 M5 膜）。

Class S 构造的第一步，是从 6d 母理论中辨认出 4d $\cN=2$ 超对称。6d 理论中总共有 16 个超荷，而一般的曲面 $\cC$ 未必保持我们想要的 4d 中的 8 个超荷；因此必须施行一次 \textbf{部分拓扑扭曲}，把一部分 R 对称性混入一部分转动对称性。让我们考虑 R 对称性的子代数 $\fso(2)_r\oplus\fso(3)_R\subset\fso(5)_R$，以及庞加莱对称性的子代数 $\fso(1,3)_{4d}\oplus\fso(2)_\cC\subset\fso(2,6)$。$\cC$ 上的扭曲转动对称性 $\fso(2)_{\mathrm{twist}}$ 定义为 $\fso(2)_\cC\oplus\fso(2)_r$ 的对角部分，于是包括 Lorentz 对称性与 R 对称性在内的残余玻色对称性为
\begin{equation}
  \fso(1,3)_{4d}\oplus\fso(2)_{\mathrm{twist}}\oplus\fso(3)_R\,.
  \tag{4.2}
\end{equation}
% [Tutor Note: 部分拓扑扭曲的记账：6d $\cN=(2,0)$ 的 R 对称性是 $\fso(5)_R$，16 个超荷在 $\fso(1,5)\oplus\fso(5)_R$ 下变换。$\cC$ 的转动代数取 $\fso(2)_\cC$，把它与 $\fso(5)_R$ 内的 $\fso(2)_r$ 作对角嵌入，就得到 (4.2) 中的 $\fso(2)_{\mathrm{twist}}$；此时超荷在 $\cC$ 上的旋量权重被 $\fso(2)_r$ 的荷抵消，16 个超荷中有 8 个成为 $\cC$ 上的标量，正是 4d 所需的 8 个超荷。$\fso(5)_R\to\fso(2)_r\oplus\fso(3)_R$ 中未被触碰的 $\fso(3)_R=\fsu(2)_R$ 就是 4d $\cN=2$ 的 $\SU(2)_R$。]
把该理论紧化在 $\cC$ 上并施行部分拓扑扭曲，我们便得到一个保持 $\fso(3)_R=\fsu(2)_R$ R 对称性、并带有两个在二重态中变换的 Weyl 超荷的系统；这恰好提供了得到 4d $\cN=2$ 庞加莱超对称代数所需的超对称数量。

6d 理论的真空由标量场 $\Phi_I$ 的 VEV 参数化，这些标量场取值于规范李代数 $\fg$ 的 Cartan 子代数，其中 $I=1,\ldots,5$。起初，这五个场可以借助 $\fso(5)_R$ 对称性被平等地相互转动；但在部分拓扑扭曲之后，情况不再如此。其中三个场，比方说 $\Phi_3$、$\Phi_4$ 与 $\Phi_5$，在未被触碰的 R 对称性 $\fso(3)_R=\fsu(2)_R$ 下带荷，因此正是这些场参数化 4d 理论的希格斯分支。其余两个场 $\Phi_1$ 与 $\Phi_2$ 在 $\fso(2)_r=\fu(1)_r$ R 对称性下带荷，所以它们参数化库仑分支。注意，在部分拓扑扭曲之后，这个 R 对称性与 $\cC$ 中的转动对称性发生了混合，也就是说这些场不再是标量。方便起见，我们记
\begin{equation}
  \Phi_z \equiv \Phi_1+\ii\Phi_2\,,
  \tag{4.3}
\end{equation}
于是 $\Phi_z\dd z$ 是 $\cC$ 上关于全纯坐标 $z$ 的 (1,0)-形式，称为 \textbf{Hitchin 场}（或
% [原文 p.51]
Higgs 场\footnote{「Higgs 场」这个名称来自关于 Higgs 丛与 Hitchin 可积系统的数学文献。由于它与物理意义上的通常 Higgs 场毫无关系（而且我们刚刚看到，它参数化的是模空间的库仑分支），在物理文献中它更多地被称为「Hitchin 场」。}）。正如我们将看到的，4d 库仑分支由 $\cC$ 上 Hitchin 场的 BPS 构型模去规范变换所决定。
% [Tutor Note: (4.3) 中的 $\Phi_z$ 是复场，它的共轭是 $\Phi_{\bar z}=\Phi_1-\ii\Phi_2$（$\Phi_1,\Phi_2$ 原本是实标量）。扭曲前 $\Phi_I$ 都是 $\cC$ 上的标量；扭曲后 $z$ 指标来自与 $\fso(2)_r$ 的混合，所以只有 $\Phi_z\dd z$（而不是 $\Phi_z$ 本身）才是 $\cC$ 上良定义的 (1,0)-形式。后文 Hitchin 方程 (4.7) 中出现的 $\bar\partial_A\Phi_z$ 与 $\partial_A\bar\Phi_z$，正是对这一几何事实的表述。]

下一步，是把描写我们理论低能动力学的 Seiberg--Witten 曲线，同 UV 构造的各个要素联系起来：Gaiotto 曲线 $\cC$、Hitchin 场 $\Phi_z$，以及规范代数——我们取 $\fg=A_n$（对其他可能李代数选择的评论推迟到 4.8.1 节）。按照 [33]（另见 [64]），让我们考虑 $T^*\cC$，即 $\cC$ 上的典范线丛。我们为 $T^*\cC$ 取坐标 $(z,x)$，其中 $x\in\C$ 参数化纤维方向。可以证明（我们将在 4.4 节进一步阐述）：Seiberg--Witten 曲线 $\Sigma\subset T^*\cC$ 是 $\cC$ 的 $n$ 重覆盖，其方程为
% [Tutor Note: 原文此处把冠词误写作 "a n-sheeted cover"（应为 "an n-sheeted"），已 900 dpi 复核底本确如此；按体例正文照译，正确形式为 an。]
\begin{equation}
  \big\langle \det\left(x-\Phi_z\right)\big\rangle = \det\left(x-\varphi_z\right) = 0\,.
  \tag{4.4}
\end{equation}
% [Tutor Note: (4.4) 左右两边不是同一个对象：左边角括号表示对场取期望值（即在库仑分支上取 VEV），右边把该 VEV 记作小写矩阵 $\varphi_z$，即 $\langle\Phi_z\rangle=\varphi_z$。$\Sigma$ 是 $n$ 重覆盖，因为在一般点上 $x$ 有 $n$ 个本征值，每个本征值在 $T^*\cC$ 的纤维上给出一个点。（注意原文的 $A_n$ 记号有滑动：§4.3 写 $\fg=A_n$ 却称 $n$ 重覆盖，而 §4.4 的 (4.8) 用标准约定 $\fg=A_{k-1}=\fsu(k)$、§4.5 又用 $A_N$ 配 $\SU(N)$；覆盖重数等于规范群基本表示的维数。）]
从 4d 的观点看，这个方程指定了一条在库仑分支上纤维化的亏格 $n$ 的曲线，正如所预期。要提取 IR 动力学所需的第二个要素是 SW 微分；在我们对 $T^*\cC$ 的坐标选择下，它简单地由下式给出
\begin{equation}
  \lambda = x\,\dd z\,.
  \tag{4.5}
\end{equation}
% [Tutor Note: (4.5) 里的 $x$ 是 $T^*\cC$ 的纤维坐标，$\dd z$ 是 $\cC$ 上的全纯微分，故 $\lambda$ 是 $T^*\cC$ 上的全纯 1-形式，沿 $\Sigma$ 拉回后就是 SW 微分；它对 $\Sigma$ 上 1-循环的积分给出 BPS 粒子的质量（或电/磁荷），见式 (4.10) 之后关于留数的讨论。]

再次地，这个构造在 M 理论中有一个漂亮的解释 [28]。在 UV，我们有一叠 $n+1$ 张包裹 $\cC$ 的 M5 膜。十一维分裂为 $4+4+3$：前 4 维是平坦的，膜在其上延展，4d 低能理论就生活在那里；接下来 4 维对应于 $\cC$ 及其横向方向，它们共同构成 $T^*\cC$；最后 3 维同样平坦，但与 M5 膜横截。标量场 $\Phi_I$ 描述 M5 膜在这些横向方向上的位移：$\Phi_3$、$\Phi_4$ 与 $\Phi_5$ 是最后三个方向上的平移，而 $\Phi_z$、$\bar\Phi_z$ 则是 $T^*\cC$ 内与 $\cC$ 垂直的两个方向上的平移。其想法是：这叠 M5 膜在进入红外时会合并成单独一张 M5 膜（我们正是这样得到低能阿贝尔 6d $\cN=(2,0)$ 理论的），后者会取一个复杂的形状，其亏格可能非平凡（当秩大于 1 时）。具体地说，M5 膜的形状——粗略说来由坐标 $x$ 支配——由 (4.4) 给出。这个方程的含义是：对于 UV 下给定的 M5 膜位移构型（即给定的 Hitchin 场 $\Phi_z$ 构型），$x$ 的可能取值就是 $\Phi_z$ 的本征值。于是，若我们考虑 $\Phi_z$ 的另一个构型，我们就会在库仑分支上移动，并在此过程中恢复我们预期的纤维化 (4.4)。

% [原文 p.52]
构造的最后一个要素是 \textbf{穿刺点}，它极大地丰富了我们在 Class S 中能够构造的理论。最容易给出其动机的视角是 M 理论，或者毋宁说是对偶的 IIA 型膜图。\footnote{在弦/M 理论图像中，Class S 构造存在一个有趣的全息对偶。对应于 Class S 理论全息对偶的 M 理论背景最早在 [34] 中被写下；T 对偶的 IIA 型的全息对偶则在 [194--196] 中得到探讨。} 在这个框架中，4d 理论由 Hanany--Witten（HW）装置得到，其中 D4 膜悬挂在 NS5 膜上。D4 膜来自包裹 M 理论圆的 M5 膜，而 NS5 膜来自在该圆上具有固定位置的 M5 膜。膜图左右两侧存在半无限 D4 膜，这会在相应箭图中导致一个非动力学的味对称性节点，因为这些 D4 膜比悬挂在两张 NS5 膜之间的 D4 膜重得多。现在，我们可以通过把该构型提升到 M 理论并对 M5 膜的内禀方向作紧化，从而把它化为更熟悉的形式；这些内禀方向来自 IIA 型膜图中沿 D4 纵向、与 NS5 横截的方向，再加上 M 理论圆。我们通过添加有限个对应于无穷远的点来完成这一紧化，每张半无限 NS5 膜或 D4 膜各对应一个点。这样，我们最终就会得到包裹在一张带若干特殊点 $p_j$ 的黎曼面上的动力学 M5 膜，这些特殊点对应于那些重的 M5 膜：它们就是穿刺点。

受此启发，我们现在可以直接在 M 理论装置中引入穿刺点。我们考虑一张带有若干标记点 $\{p_i\}_i$ 的黎曼面 $\cC$，并在原先那叠 M5 膜延展于 $\R^{1,3}\times\cC$ 的构型中加入延展于 $\R^{1,3}\times T^*_{p_i}\cC$ 的缺陷 M5 膜，以此实现穿刺点。更抽象地说，我们也可以把 6d $\cN=(2,0)$ SCFT 中的穿刺点一般地理解为带指定边界条件的余维二缺陷。它们将对应于 Hitchin 场在点 $p_i\in\cC$ 处的极点，这些极点携带关于 4d 理论整体对称性的信息。粗略的想法是：借助 (4.4)，它们会成为 $x(z)$ 的极点；当我们绕 $p_i$ 积分 Seiberg--Witten 微分时，就会拾取它们的留数，而该留数因此被认同为 BPS 粒子的质量参数。可以存在若干类对应于不同味对称性的穿刺点，但我们把更详细的讨论推迟到 4.5 节。

% （§4.4 标题与「与此同时…」段仍在印刷 p.52，此处不换页）
\SUB{4.4}{从经典 Hitchin 系统得到量子几何}

与此同时，我们想更精确地说明：如何从紧化在 UV 曲线 $\cC$ 上的 6d $\cN=(2,0)$ 理论中，恢复对 4d 描述有用的有效理论。这要借助 Hitchin 可积系统 [197] 来完成：由此我们可以通过求解一些经典方程，得到精确的库仑分支几何。为了解释这一关系，我们利用图 7 所示的一连串紧化，从 6d $\cN=(2,0)$ 理论一路到 3d $\cN=4$ 理论。

% [原文 p.53]
\begin{center}
  \includegraphics[width=0.55\linewidth]{fig7.png}
\end{center}
\figtitle{7}{6d $\cN=(2,0)$ SCFT、5d $\cN=2$ SYM、4d $\cN=2$ Class S 理论及其 3d $\cN=4$ 对应理论之间通过相继紧化建立的关系。镜像对称联系着同一个 3d $\cN=4$ SCFT 在红外下的两种 UV 描述。}

% [Tutor Note: 图 7 的四条边：上左、上右两条箭头分别表示在 $\cC$ 上和在 $S^1$ 上紧化（箭头旁标注的正是 $\cC$ 与 $S^1$）；下方左右两条弧线表示把 4d 理论再紧化到 $S^1$ 上、把 5d 理论再紧化到 $\cC$ 上，两者落到同一个 3d $\cN=4$ 理论；中间的双向箭头是 3d 镜像对称，它交换两边的 Coulomb 分支与 Higgs 分支，因而把「用 4d 线算子描述的 HK 流形」与「5d SYM 在 $\cC$ 上的 Hitchin 模空间」等同起来。]
这些箭头可以如下理解。把 6d $\cN=(2,0)$ 理论紧化在黎曼面 $\cC$ 上，会得到我们在前几节中介绍过的一大类强耦合 4d $\cN=2$ 理论。为了从 6d 装置中提取所得 4d 理论的物理信息，事实证明进一步紧化到半径为 $R$ 的圆 $S^1$ 上是有用的。3d $\cN=4$ 理论一般是一个非平凡的 SCFT，它编码 4d 母理论的物理。特别地，在 3d 极限下，4d 理论的库仑分支会因 $S^1$ 上 BPS 线算子的 VEV 而得到增强；而完整的 3d CB EFT 由 $\cN=4$ sigma 模型描述，其目标空间是复维数为 $2r$（即 4d CB 维数的两倍）的超 Kähler（HK）流形 $\cM$ [35]。从 3d 的角度看，$\cM$ 是真空模空间中保持 (4.2) 里 $\fso(3)_R$ 对称性的那个分支。

用 4d 线算子对 HK 流形的这种描述可能有点抽象。事实上，如果我们把紧化的次序颠倒过来，HK 流形 $\cM$ 的结构会变得更透明。\footnote{注意，从六维到三维的紧化次序对模空间 $\cM$ 无关紧要 [35]。这与如下事实相联系：由于拓扑扭曲，BPS 保护量不关心 $\cC$ 与 $S^1$ 之间的相对标度。} 如前所述，6d $\cN=(2,0)$ 理论是非拉格朗日量的，但它在半径为 $R$ 的圆 $S^1$ 上的紧化，在低能下由 5d 最大超对称 Yang--Mills 理论描述。
% [接续] 本段跨 tr05a/tr05b 两个文件：上一句起于原文 p.54 末、止于 p.55 首行（句界跨页，按「开始者写完」已在此写完整）；下一句起于 p.55，由 tr05b 无空行续写，仍属同一自然段。
